% Chapter 19: The Sensitivity Conjecture
\let\LocalMacro\relax

\chapter{Sensitivity Conjecture}

\section{The Problem}

\subsection{Informal}
Imagine you are asked a yes-or-no question about a long string of coin flips---say, ``is the number of heads even?'' If you flip one coin, the answer might change or it might not. One way to measure the question's complexity is the most coins you could flip at once to force the answer to change. Another way measures the same thing but lets you flip whole groups of coins together. The sensitivity conjecture, open for 27 years, asked whether these two ways of measuring complexity are always closely related: if a question looks simple by one measure, must it look simple by the other too?

\subsection{Formal}
\begin{conjecture}[Nisan-Szegedy, 1992]
$s(f)$ and $bs(f)$ are polynomially related: $bs(f) \le s(f)^c$ for some constant $c$.
\end{conjecture}

\subsection{Historical Context}
The sensitivity conjecture stood for 27 years, resisting attacks from Fourier analysis, probability theory, and complexity theory. It was one of the oldest open problems in theoretical computer science.

\subsection{What Was Known}
The conjecture was posed by Nisan and Szegedy in 1992. For decades, attempts to prove it used Fourier analysis, probability theory, and algebraic approaches---all without success. The best known bound was exponentially weaker than the conjectured polynomial relationship. The problem was considered one of the most important open questions in complexity theory.

\subsection{Why It Matters}
The sensitivity conjecture is about measuring the complexity of yes-or-no questions---a fundamental concept in computer science. If two different ways of measuring complexity are always closely related, it means that the complexity of a problem is a robust notion, not an artifact of how you choose to measure it. This has implications for algorithm design and the theory of computation.

\subsection{Simplified Version}
For the parity function $f(x_1, \ldots, x_n) = x_1 \oplus \cdots \oplus x_n$ (which outputs 1 if an odd number of bits are 1), the sensitivity is $\sqrt{n}$ while the block sensitivity is $n$---a quadratic gap. The conjecture asserted that this gap is always at most polynomial: $bs(f) \le s(f)^c$ for some constant $c$. Huang's 2019 proof showed $c \le 4$, settling the question.

\section{Mathematical Theory}

\subsection{Prerequisites}
This chapter assumes familiarity with:
\begin{itemize}
\item Boolean functions ($f: \{0,1\}^n \to \{0,1\}$)
\item Spectral graph theory (eigenvalues of adjacency matrices)
\item Basic complexity theory (decision trees, polynomial degree)
\end{itemize}

\subsection{Boolean Function Complexity Measures}

The sensitivity conjecture concerns the relationship between several complexity measures of Boolean functions. For $f: \{0,1\}^n \to \{0,1\}$, the key measures are:

\begin{definition}[Sensitivity]
The \emph{sensitivity} $s(f)$ is the maximum, over all $x \in \{0,1\}^n$, of the number of coordinates $i$ such that flipping the $i$-th bit of $x$ changes the value of $f$.
\end{definition}

\begin{definition}[Block Sensitivity]
The \emph{block sensitivity} $bs(f)$ is the maximum number $k$ such that there exist disjoint non-empty sets of coordinates $B_1, \ldots, B_k$ with the property that for each $i$, flipping all bits in $B_i$ changes the value of $f$.
\end{definition}

\begin{definition}[Degree]
The \emph{degree} $\deg(f)$ is the degree of the unique multilinear polynomial $P$ over $\mathbb{R}$ such that $f(x) = P(x)$ for all $x \in \{0,1\}^n$.
\end{definition}

\begin{definition}[Decision Tree Depth]
The \emph{decision tree depth} $D(f)$ is the minimum depth of a deterministic decision tree that computes $f$.
\end{definition}

These measures are ordered: $s(f) \le bs(f)$, $s(f) \le \deg(f)$, $s(f) \le D(f)$. The conjecture asserted that all are polynomially related.

\subsection{The Hypercube and Its Spectrum}

The $n$-dimensional hypercube graph $Q_n$ has vertex set $\{0,1\}^n$ with edges between strings differing in exactly one coordinate. Its spectral theory is classical:

\begin{fact}
The eigenvalues of the adjacency matrix of $Q_n$ are $n - 2k$ for $k = 0, 1, \ldots, n$, with multiplicity $\binom{n}{k}$.
\end{fact}

The eigenspace for eigenvalue $n$ is spanned by the all-ones vector. The eigenspace for eigenvalue $n - 2$ is spanned by the $n$ coordinate functions $x_i$.

\subsection{Cauchy Interlacing}

The key tool in Huang's proof is a classical result about how eigenvalues change when restricting to a subgraph:

\begin{theorem}[Cauchy Interlacing]
If $A$ is an $n \times n$ symmetric matrix with eigenvalues $\lambda_1 \ge \cdots \ge \lambda_n$, and $B$ is an $(n-1) \times (n-1)$ principal submatrix with eigenvalues $\mu_1 \ge \cdots \ge \mu_{n-1}$, then:
\[
\lambda_1 \ge \mu_1 \ge \lambda_2 \ge \mu_2 \ge \cdots \ge \mu_{n-1} \ge \lambda_n.
\]
\end{theorem}

This means that restricting to a subgraph can only push eigenvalues inward, but the largest eigenvalue can decrease by at most the gap between the two largest eigenvalues.

\subsection{Signings and Hadamard Matrices}

A \emph{signing} of a graph $G$ is an assignment of $\pm 1$ to each edge, producing a signed adjacency matrix $A_G$. Huang's key construction produces a signing of $Q_n$ such that:

\begin{fact}
There exists a signing of $Q_n$ such that $A^2 = nI$, making $A$ proportional to a Hadamard matrix.
\end{fact}

This is the crucial ingredient: it means the signed hypercube has all eigenvalues equal to $\pm\sqrt{n}$, giving maximum spectral gap.

\section{The Solution}

\subsection{Huang (2019)}

\begin{theorem}[Huang \cite{huang2019}]
$s(f) \ge \sqrt{\deg(f)}$, hence $bs(f) \le s(f)^4$.
\end{theorem}

Huang's proof is famously short (2 pages) and uses only spectral graph theory:

\begin{enumerate}
\item Consider the $n$-dimensional hypercube graph $Q_n$, whose vertices are $\{0,1\}^n$ and edges connect strings differing in one coordinate.
\item Construct a signing of the adjacency matrix $A$ of $Q_n$ such that $A^2 = nI$ (a Hadamard matrix structure).
\item By the Cauchy interlacing theorem, the subgraph induced by $f^{-1}(1)$ has an eigenvalue $\ge \sqrt{n}$.
\item This implies the subgraph has minimum degree $\ge \sqrt{n}$, so $s(f) \ge \sqrt{n}$.
\item Since $\deg(f) = n$ for the dictator function, $s(f) \ge \sqrt{\deg(f)}$.
\end{enumerate}

\begin{highlightbox}
Huang's proof resolved a 27-year-old problem using a single clever idea: a signing of the hypercube's adjacency matrix. The proof uses no Fourier analysis, no probability, and no complexity theory---just linear algebra.
\end{highlightbox}

\section{Impact}
\begin{itemize}
\item \textbf{Complexity theory}: Established polynomial relationships between all known boolean function complexity measures.
\item \textbf{Spectral methods}: The signing technique has been applied to other graph-theoretic problems.
\item \textbf{Open question}: The exact polynomial relating $s(f)$ and $bs(f)$ remains unknown---Huang's proof gives $c = 4$, but the conjectured value is $c = 1$.
\end{itemize}

\section{Exercises}

\begin{exercise}[Basic]
Compute $s(f)$ and $bs(f)$ for the parity function $f(x_1, \ldots, x_n) = x_1 \oplus \cdots \oplus x_n$.
\end{exercise}

\begin{exercise}[Intermediate]
Explain the Cauchy interlacing theorem: if $A$ is an $n \times n$ symmetric matrix and $B$ is an $(n-1) \times (n-1)$ principal submatrix, how do their eigenvalues relate?
\end{exercise}

\begin{exercise}[Challenge]
Sketch Huang's construction of the signing of $Q_n$'s adjacency matrix. Why does $A^2 = nI$ hold?
\end{exercise}

\section{Timeline}

\begin{tabularx}{\linewidth}{lX}
\toprule
\textbf{Year} & \textbf{Event} \\ \midrule
1992 & Nisan-Szegedy formulate the sensitivity conjecture \\
2019 & Huang proves $s(f) \ge \sqrt{\deg(f)}$ in 2 pages \\
\bottomrule
\end{tabularx}

\section{Citations}

\begin{enumerate}
\item Huang, H. (2019). ``Induced subgraphs of hypercubes and a proof of the sensitivity conjecture.'' Annals of Mathematics, 190(3), 949--955.
\item Nisan, N., \& Szegedy, M. (1992). ``On the degree of a boolean function as a real polynomial.'' Computational Complexity, 4(4), 301--313.
\end{enumerate}
